% Part: sets-functions-relations
% Chapter: sets
% Section: subsets

\documentclass[../../../include/open-logic-section]{subfiles}

\begin{document}

\olfileid[id]{sfr}{set}{sub}
\olsection{Himpunan Bagian dan Himpunan Kuasa}

\begin{explain}
Kita akan sering membandingkan himpunan. Salah satu perbandingan yang
wajar adalah sebagai berikut: \emph{semua yang ada di dalam satu
himpunan juga ada di dalam himpunan yang lain}. Situasi ini cukup
penting sehingga kita memperkenalkan notasi baru.
\end{explain}

\begin{defn}[Himpunan bagian]
Jika setiap !!{element} himpunan $A$ juga merupakan !!a{element}
dari~$B$, kita mengatakan bahwa $A$ adalah \emph{himpunan bagian}
dari~$B$ dan menulis $A \subseteq B$. Jika $A$ bukan himpunan bagian
dari~$B$, kita menulis $A \not\subseteq B$. Jika $A \subseteq B$
tetapi $A \neq B$, kita menulis $A \subsetneq B$ dan mengatakan bahwa
$A$ adalah \emph{himpunan bagian sejati} dari $B$.
\end{defn}

\begin{ex}
Setiap himpunan merupakan himpunan bagian dari dirinya sendiri, dan
$\emptyset$ merupakan himpunan bagian dari setiap himpunan. Himpunan
bilangan asli genap merupakan himpunan bagian dari himpunan bilangan
asli. Selain itu, $\{ a, b \} \subseteq \{ a, b, c \}$. Akan tetapi,
$\{ a, b, e \}$ bukan himpunan bagian dari $\{ a, b, c \}$.
\end{ex}

\begin{ex}
Bilangan $2$ merupakan !!{element} dari himpunan bilangan bulat,
sedangkan himpunan bilangan genap merupakan himpunan bagian dari
himpunan bilangan bulat. Meskipun demikian, sebuah himpunan dapat
\emph{sekaligus} merupakan !!a{element} dan himpunan bagian dari
himpunan lain. Sebagai contoh, $\{0\} \in \{0, \{0\}\}$ dan juga
$\{0\} \subseteq \{0, \{0\}\}$.
\end{ex}

Ekstensionalitas memberikan kriteria identitas bagi himpunan: $A = B$
jika dan hanya jika setiap !!{element} dari~$A$ juga merupakan
!!a{element} dari~$B$, dan sebaliknya. Definisi ``himpunan bagian''
mendefinisikan $A \subseteq B$ tepat sebagai separuh pertama kriteria
ini: setiap !!{element} dari~$A$ juga merupakan !!a{element} dari~$B$.
Tentu saja definisi itu juga berlaku jika kita menukar $A$ dan $B$:
$B \subseteq A$ jika dan hanya jika setiap !!{element} dari~$B$ juga
merupakan !!a{element} dari~$A$. Pernyataan terakhir ini tepat
merupakan bagian ``sebaliknya'' dari ekstensionalitas. Dengan kata
lain, ekstensionalitas menyatakan bahwa dua himpunan sama jika dan
hanya jika keduanya merupakan himpunan bagian satu sama lain.

\begin{prop}
$A = B$ jika dan hanya jika $A \subseteq B$ dan $B \subseteq A$.
\end{prop}

Sekarang juga merupakan saat yang tepat untuk memperkenalkan beberapa
notasi lain yang berguna. Dalam mendefinisikan kapan $A$ merupakan
himpunan bagian dari~$B$, kita mengatakan bahwa ``setiap !!{element}
dari~$A$ adalah \dots,'' lalu mengisi ``$\dots$'' dengan
``!!a{element} dari $B$''. \emph{Bentuk} ungkapan ini begitu umum sehingga
berguna untuk memperkenalkan notasi formal baginya.

\begin{defn}\ollabel{forallxina}
$(\forall x \in A)\phi$ menyingkat $\forall x(x \in A \lif
\phi)$. Demikian pula, $(\exists x \in A)\phi$ menyingkat
$\exists x(x \in A \land \phi)$.
\end{defn}

Dengan notasi ini, kita dapat mengatakan bahwa $A \subseteq B$ jika
dan hanya jika $(\forall x \in A)x \in B$.

Sekarang kita beralih kepada suatu jenis himpunan tertentu: himpunan
semua himpunan bagian dari suatu himpunan yang diberikan.

\begin{defn}[Himpunan kuasa]
Himpunan yang terdiri atas semua himpunan bagian dari himpunan~$A$
disebut \emph{himpunan kuasa dari}~$A$ dan ditulis $\Pow{A}$.
  \[
    \Pow{A} = \Setabs{B}{B \subseteq A}
  \]
\end{defn}

\begin{ex}
Apa saja himpunan bagian yang mungkin dari $\{ a, b, c \}$? Semuanya
adalah $\emptyset$, $\{a \}$, $\{b\}$, $\{c\}$, $\{a, b\}$,
$\{a, c\}$, $\{b, c\}$, dan $\{a, b, c\}$. Himpunan yang terdiri
atas semua himpunan bagian ini adalah $\Pow{\{a,b,c\}}$:
\[
\Pow{\{ a, b, c \}} = \{\emptyset, \{a \}, \{b\}, \{c\}, \{a, b\},
\{b, c\}, \{a, c\}, \{a, b, c\}\}
\]
\end{ex}

\begin{prob}
Tuliskan semua himpunan bagian dari $\{a, b, c, d\}$.
\end{prob}

\begin{prob}
Tunjukkan bahwa jika $A$ mempunyai $n$ !!{element}s, maka $\Pow{A}$
mempunyai $2^n$ !!{element}s.
\end{prob}

\end{document}
